\documentclass{article}
\usepackage{lmodern, amssymb,amsmath, graphicx, hyperref}
% === bibliography package ===
\usepackage{natbib}
% \usepackage[colorinlistoftodos, prependcaption]{todonotes} % to use \todo 
\usepackage[margin=1in]{geometry}
\usepackage{csquotes}
\usepackage{setspace}
\hypersetup{
    colorlinks=true,
    linkcolor=blue,
    filecolor=magenta,      
    urlcolor=cyan,
    citecolor = black
}
\urlstyle{same}  % don't use monospace font for urls

  \title{Project Summary}
    %\author{Devin Judge-Lord\thanks{Ph.D. Candidate, University of Wisconsin-Madison}\and Justin Grimmer\thanks{Professor, Stanford University} \and Eleanor Neff Powell\thanks{Booth Fowler Associate Professor, University of Wisconsin-Madison}}
    \date{}
    
\graphicspath{{gh-pages/Figs/}}


\begin{document}

\maketitle


%%%  4600 Character Max with Spaces
%%% Must have following headers: Overview, Intellectual Merit, Broader Impacts

%\begin{abstract}
%\noindent Representation in Congress is based, in part, on legislators' ability to assist their constituents with the federal bureaucracy.  Constituency service is central to definitions of representation as well as empirical explanations for phenomena such a the incumbency advantage. And yet, too little is known about when and how elected officials contact the bureaucracy on behalf of their constituents. We assemble a massive new database of over 150,000 Congressional requests to agencies between 2007 and 2017 obtained through over 250 FOIA requests, approaching a census of all contacts with the bureaucracy. We examine when legislators contact agencies, which legislators contact those agencies, and the purpose of the contact.  Our work provides new insights into who provides constituency service, why they provide that service, and the consequences of the agency requests for how governments make decisions. % Not only do representative vary significantly in their productivity, the top advocates for individual constituents are also the top advocates for policy change, making contact with the bureaucracy an important dimension of representation.
%\end{abstract}

%\newpage
\doublespacing

\section*{Overview}

Representation in Congress is based, in part, on legislators' ability to assist their constituents with the federal bureaucracy. Constituency service is central to definitions of representation as well as empirical explanations for phenomena such as the incumbency advantage. While more visible forms of representation (e.g. roll-call voting) have received considerable scholarly attention, too little is known about when and how elected officials contact the bureaucracy on behalf of their constituents. We are seeking funding to address this imbalance by assembling a massive new database of Congressional communication with federal agencies--in essence creating the first census of all contacts with the bureaucracy over a ten-year period. We examine when legislators contact agencies, which legislators contact those agencies, and the purpose of the contact. Our work will provide new insights into three areas: 1) inequalities in representation, 2) money in politics and 3) legislative behavior. Our preliminary findings from our pilot study comprising the largest dataset of congressional communication with federal agencies to date (over 150,000 requests) has revealed some surprising results -- staggering inequality in the representation legislators are providing their constituents.   



% Not only do representative vary significantly in their productivity, the top advocates for individual constituents are also the top advocates for policy change, making contact with the bureaucracy an important dimension of representation.

\section*{Intellectual Merit}

Congressional scholars have written frequently on forms of congressional representation in general and constituency service in particular. Indeed, it is one of the most frequently discussed forms of congressional representation in the literature, and yet a mere handful of these analyses are based on constituency service data. This omission in the literature is for a good-reason. Up to this point, we've not had any data on constituency service. Our approach solves this problem by turning to the federal agencies themselves, and asking them to tell us what contact they receive from members of Congress. We will use this data to provide insights in three fields: 1) inequalities in representation, 2) money in politics and 3) legislative behavior. 

First, constituency service data has the potential to provide a new approach to measuring and assessing inequalities in representation. Unlike high profile roll call votes, constituency service is largely hidden from public view and there is almost no public accountability. Our preliminary results from a pilot study suggest dramatic results: inequality in representation is comparable to the high level of income inequality in the US. These results suggest some citizens are receiving high quality representation, while others are largely ignored by their representatives. 

Second, in addition to exploring questions of representational inequality, we can answer long-standing questions in the field of money in politics. These records reveal when legislators contact federal agencies on behalf of corporate contributors. By matching up communication data with campaign contributions, we can observe which members advocate on behalf of their donors, and what effect that advocacy may have. 

Third, this data can provide insights into legislator behavior.  There have been long-standing public concerns that incumbent legislators ``go Washington'' and may neglect their constituents over time. This data allows us to explicitly examine that claim, and results from our pilot study suggest the opposite – as legislators gain influence in Washington they increase their advocacy on behalf of their constituents. 

\section*{Broader Impacts}

In addition to the scholarly value of this data collection project both for our own research and for the scholarly community writ large, we will make it easily accessible for both journalists and the general public. To make the data accessible, we will not simply post the bulk data online, but we'll create an App to make it easy for any citizen to look up their own member of Congress and observe what constituency service he or she has provided, as well as how they compare to other members. This unpacking of the data and making it easily searchable will also facilitate investigative journalism into these three topics of inequality of representation, money in politics and legislative behavior. Historically, journalists have sought out records of what constituency service a legislator is providing by placing an individual FOIA request. However, that process is slow, cumbersome and unwieldly given the speed of modern journalism. This data, website, and app will revolutionize that process and provide the public with new insight into how they are represented. 

\bibliographystyle{chicago-apsr}

\bibliography{congress2018}

%\newpage
%\listoftodos[Notes]


\end{document}






